\chapter{Completion of the strong C2 ordinary--small source chart}
\label{component:p6b22e944bd7c}
\noindent\textbf{Component }\path{p6b22e944bd7c}\par
{\small\noindent Source: \path{R13/baseline/baseline/advances/ui/ui_all_tests_recovered.tex}\par}
{\small\noindent SHA256: \path{6b22e944bd7c51d3c7e80d46a5a2db0b3ae5a6f6171dfe98a55787befe99a6e4}\par}
{\small\noindent Inventory use: UA1--UA3 only: elementary real rank-one allocation lemma and proof | UA1--UA16 complete source openings, finite-cap approximation, remote/testing/projective endpoints and all state ratios | C1\_\allowbreak{}BARE\_\allowbreak{}OTHER\_\allowbreak{}TERMINALS; precise subordinate uses retained below\par}
\medskip
\section{Completion of the strong C2 ordinary--small source chart}
\label{p6b22e944bd7c:sec:ui-all-tests-recovered}

\begin{theorem}[Strong C2 UI, all independent tests]
Use the original C2 source (UR1)--(UR3), with $d=1$, actual source
angles $u_c=a>0$ a bounded unit and $v_c=k>0$ infinitesimal, and
$C,\alpha$ strong. Then the entire parent (UR6) is a sum of squares
in the tensor product of the independent source and testing fields,
for every testing angle, including projective endpoints, and every
homogeneous state. The statement retains the literal boundary $n=0$
and arbitrary positive rates of every source opening.
\end{theorem}

The local contact certificate in Section~\ref{p76545700d90b:sec:local-ui-certificate}
already covers $v$ infinitesimal and $u$ with the same residue as $a$.
We supply its complete complement. No comparison between $v$ and $k$
is made. The argument uses one elementary rank-one allocation, followed
by short positive frames at the testing endpoints.

\paragraph{A rank-one allocation lemma.}
Let $P$ be a positive definite real symmetric matrix, let $\ell,b$ be
real columns, and let $H,c_0>0$. Set
\[
 d=\ell^TP^{-1}\ell,\qquad e=\ell^TP^{-1}b,\qquad
 h=c_2-b^TP^{-1}b.
\]
Suppose
\[
 p(T)=c_0+(h+4Hdc_0)T+4H(dh+e^2)T^2>0\quad(T\ge0),
 \qquad dh+e^2>0.                                      \tag{UA1}
\]
There exists a real $\lambda$ such that
\[
 |\lambda|<\sqrt{4Hc_0},\qquad
 h+2e\lambda-d\lambda^2>0.                             \tag{UA2}
\]
Indeed put $s=\sqrt{4Hc_0}$. If $d>0$ and $|e|<ds$, take
$\lambda=e/d$. If $|e|\ge ds$, the value at the appropriate endpoint
is $h+2|e|s-ds^2$. Were this nonpositive, writing
$x=|e|-ds\ge0$ would give
$h+ds^2\le-2sx$ and $dh+e^2\le x^2$; the quadratic in (UA1)
would then have a nonpositive value at a positive argument. Thus a
nearby interior $\lambda$ works. This includes equality $|e|=ds$.
The scalar interval argument also holds in a limiting case $d=0$;
then $e\ne0$ follows from $dh+e^2>0$. (For a finite positive $P$
itself, $d=0$ would force $\ell=e=0$, so this case occurs only
when using the continuous scalar extension at a boundary.)

Consequently the two real matrices
\[
 \begin{pmatrix}P&b-\lambda\ell\\(b-\lambda\ell)^T&c_2\end{pmatrix},
 \qquad
 \begin{pmatrix}4H&\lambda\\\lambda&c_0\end{pmatrix}       \tag{UA3}
\]
are positive definite. This is an elementary real lemma. Its inverse
and square root are never taken in the tensor ring.

\paragraph{The finite ordinary angle away from contact.}
Assume $v$ infinitesimal and $u$ a bounded unit whose real residue has
square different from $a_0^2$. Keep the centered packets
$P_L,P_C,P_n,P_m,P_\alpha$ of (UR12), (UR14) outside the Gram, and
combine $P_H$ with the exact remainder. Use the six separated roots
\[
 r_0=uW-Y,\quad X,\quad Z,\quad V_0=vW,\quad K_0=kW,
 \quad R_0=\frac{v(uX-Z)}{k}.                           \tag{UA4}
\]
They satisfy the exact zero relation
\[
 R_0K_0-(uX-Z)V_0=0.                                  \tag{UA5}
\]
All divisions in (UA4) are pure source divisions.

For completeness the boundedness of the $C,D$ remainder in this frame
is explicit. Write $\ell_0=L\tau$, $c_R=C\rho k^2$, and
$Q=uX-Z$. Then (UR13) gives exactly
\begin{align*}
 G_{CD}+P_H={}&4\gamma R_0^2
 +8\ell_0kK_0r_0+8\ell_0qk^2V_0Q-4\ell_0\tau k^2K_0^2\\
 &+8c_RvV_0r_0-4c_R\rho k^2V_0^2.                  \tag{UA6}
\end{align*}
All displayed coefficients are bounded by (UR9)--(UR11); all terms
except the first have zero residue. The $A$ packet is exactly
$4(A/k^2)D_uD_v(uK_0-kr_0)^2$. Substituting
$Y=uW-r_0$ in (UR15) and the signed core gives a bounded Gram on
(UA4), with no discarded normal term. Its residue has an isolated
$2r_0^2$ term. First suppose the leading source weights are finite.
The remaining residue is
\[
 q_s^TPq_s+2V_0b_*^Tq_s+c_2V_0^2+4H R_0^2+c_0K_0^2,
 \quad q_s=(X,Z)^T,                                  \tag{UA7}
\]
where $D_u=1+u^2$, $H=H_0$, $A_w=A_0^{\rm weak}$ are (RC3), and
\[
\begin{gathered}
 P=\begin{pmatrix}
 4nD_u+4m+2&(4m+2)u\\
 (4m+2)u&(4m+2)u^2+4\alpha D_u
 \end{pmatrix},\quad \ell=(u,-1)^T,\\
 b_*=\binom{-4mD_u}{-D_uu(4m+4\alpha+2)},\quad
 c_0=4A_wD_uu^2,\quad
 c_2=4mD_u^2+(4\alpha+2)D_uu^2.                      
\end{gathered}
\tag{UA8}
\]
The matrix $P$ is strictly positive because $m,\alpha$ are strong.
The last sign-changing part of (UA7) is handled by adding
$2\lambda$ times (UA5). This gives precisely the two blocks (UA3).

Here is a direct proof of (UA1) for every ordinary angle away from
contact. Put $N_*=2n+1$, $T_*=2m+2\alpha+1$, $U=u^2$, and
$T=V^2$. The leading state polynomial (RC4), with its testing $a$
replaced by $u$ but its source unchanged, has determinant
\[
 \det f=8D_u^2u^2V^2\,\mathcal L(U,T),                \tag{UA9}
\]
where
\begin{align*}
\mathcal L={}&\frac{4\alpha mN_*}{U}
 +4A_wHT_*U+\frac{4A_w\alpha(2m+N_*)}{T}
 +\frac{4A_wnT_*U}{T}\\
 &+2HN_*T_*T+\frac{4HmN_*T}{U}
 +4A_wH(2m+N_*)-N_*T_*.
\end{align*}
After division by $N_*T_*$ all possible opening poles cancel. Define
\[
 i_N=N_*^{-1},\quad t_m=m/T_*,\quad t_{am}=\alpha m/T_*,
 \quad q_\alpha=\frac{\alpha(2m+N_*)}{N_*T_*},\quad i_\alpha=1/\alpha.
\]
Writing $z=a^2$, $A_c=1+zr$, and using the leading $\Delta$ of (RC3),
these quantities are exactly
\[
\begin{aligned}
 i_N&=\frac{r\Delta}{(b+r)A_c^2},&
 t_m&=\frac{(1-r)(1-zb)}{2b(b+r)(z+1)},\\
 t_{am}&=\frac{z^2(b+r)(1-r)}{4\Delta},&
 q_\alpha&=\frac{bz^2\{(z+1)b+(z+2)r-1\}}{2A_c^2},\\
 i_\alpha&=\frac{2(1-zb)\Delta}{z^2(z+1)b(b+r)^2}.
\end{aligned}\tag{UA10}
\]
All denominators here are units. The normalized Laurent polynomial is
\[
 \frac{\mathcal L}{N_*T_*}
 =\frac{4t_{am}}U+4A_wHi_NU+\frac{4A_wq_\alpha}T
  +2A_w(1-i_N)\frac UT+2HT+4Ht_m\frac TU
  +4A_wHq_\alpha i_\alpha-1.                         \tag{UA11}
\]
Its value and its two first derivatives vanish at $(U,T)=(z,1)$,
as direct substitution verifies. In logarithmic coordinates the
nonconstant terms have nonnegative coefficients. The coefficients of
$U^{-1}$ and $T$ are positive units. If $i_N>0$, the coefficients of
$U$ and $T^{-1}$ are positive; if $i_N=0$, that of $U/T$ is positive.
The function is therefore strictly convex in logarithmic coordinates
and has its unique minimum at $(z,1)$. This proves strict positivity
at fixed $U\ne z$ for every $T>0$. The constant and leading
coefficients of the equivalent quadratic (UA1) are strictly positive:
$c_0>0$, and the positive $T$ coefficient in (UA11) gives the second
assertion. Thus (UA2)--(UA3) prove the required strict residue Gram.

\paragraph{Opening boundaries and divergent source weights.}
The preceding proof extends to all allowed opening residues, rather
than assuming $n,m,\alpha$ bounded. The following formulas make the
needed uniformity checkable. Put
\[
 J_0=q_\alpha+\frac{u^2(1-i_N)}2.
\]
Then $J_0$ is a positive unit, and direct two-by-two inversion gives
\[
\begin{aligned}
 \frac{\det P}{N_*T_*}&=8D_uJ_0,\\
 d&=\frac{u^2i_N+q_\alpha i_\alpha}{4J_0},&
 e&=\frac{uD_u}{2J_0},&
 h&=\frac{D_u(4t_{am}-u^2)}{J_0},\\
 (P^{-1}b_*)_1&=\frac{i_N(u^2-4t_{am})}{2J_0},&
 (P^{-1}b_*)_2&=-\frac{u\{4t_{am}i_N+(1-i_N)u^2+1\}}{2J_0}.
\end{aligned}\tag{UA12}
\]
To check the unit assertion, if $i_N=0$ then $J_0=u^2/2>0$.
If $i_N>0$, the weight $N_*$ is finite and the strong lower bounds
on $m,\alpha$ give $q_\alpha>0$: this follows directly from
$\alpha(2m+N_*)/[N_*(2m+2\alpha+1)]$, whether both weights are
finite or one diverges. Thus (UA12) has only unit denominators on
the entire source residue domain.

The strictly positive leading coefficient is also explicit:
\[
 dh+e^2=\frac{D_u(u^2+2t_m)}{2J_0}>0.
\]
Equations (UA10)--(UA12) prove continuity of $d,e,h$ and of the
minimizing hard vector. The convexity proof for (UA11) remains strict
at every allowed boundary. Choose $\lambda$ from the strict interval
(UA2) at that boundary; the same strict inequalities hold nearby.
Moreover $P$ is bounded below by a positive real matrix, while
$P^{-1}b_*$ and $P^{-1}\ell$ are bounded. Its Schur factorization
therefore bounds the first block in (UA3) below by a positive real
multiple of the identity, uniformly in the nearby source. The second
block is a strictly positive real matrix as well.

The finite leading sources used in this limiting argument are valid
real interior sources. This point must be established separately:
the actual source has $b=h+(k^2/z)t$, and it is not legitimate to
substitute that actual $b$ into a leading formula with denominator
$1-zb$ when this denominator has an infinitesimal sign change.
For the limiting real parameters, put
\[
 f(b,r)=1+r(z-1)-brz(z+1)-r^2z(z+2).
\]
Their domain is $b>0$, $0\le r<1$, $g=(1+z)b-1+r\ge0$,
$1-zb\ge0$, and $f(b,r)\ge0$. For $0<s<1$, let
\[
 b_s=sb+(1-s)/(1+z),\qquad r_s=sr.
\]
Direct expansion gives
\[
 g_s=sg,\qquad 1-zb_s>0,\qquad
 f(b_s,r_s)=(1-s)(1+s-sr)+s^2f(b,r)>0.
\]
Increase $b_s$ by a sufficiently small $\epsilon>0$. This makes
$g_s$ strict, decreases the last two strict margins by respectively
$z\epsilon$ and $r_sz(z+1)\epsilon$, and therefore preserves them.
If $r=0$, then increase $r_s$ by a sufficiently small positive
amount; $g$ increases and $f$ remains positive by continuity.
Letting these perturbations tend to zero as $s\to1$ supplies valid
finite real leading sources approaching every allowed residue.
Their $t$ may be chosen as $(1+zr)/2$, since the leading harmonic
weight is independent of $t$. The normalized leading identity
$2rn\Delta=bf(b,r)$ verifies the nonnegative-$n$ requirement.
The same construction supplies the real-interior approximation
used in the finite-cap step of the local UI certificate.

At most two of $n,m,\alpha$ diverge: $m$ can diverge only at $g=0$,
and $\alpha$ only at $1-zb=0$, which cannot occur together. Their
centered excess rows have hard-state parts proportional to
\[
 X,\qquad X+uZ,\qquad Z.
\]
Any possible pair is independent because $u$ is a unit. For a vector
in their common limiting kernel, change only $(X,Z)$ by a vanishing
amount to annihilate the corresponding nearby exact rows. Their
large penalties then vanish exactly. Taking the limit in the uniform
strict lower bound proves that the bounded finite part is strictly
positive on that limiting kernel. The elementary finite-dimensional
penalty lemma used in the local theorem now supplies finite real caps
$n_*,m_*,\alpha_*$. Their centered excesses are precisely the retained
packets (UR14). The exact bounded remainder with those caps, after
adding (UA5), has a strictly positive residue Gram on (UA4).
This proves the finite ordinary-angle complement for all source
openings and literal $n=0$.

\paragraph{Endpoint frames.}
Choose positive real $\mu,m_*,\alpha_*$ below, respectively,
$k^2C,k^2D$, $m$, and $\alpha$. The products $k^2C,k^2D$ are strong
by (UR5), and $A_w=A/k^2$ is a positive unit. Retain $n$ and all
positive excesses in full. The following formulas concern the parent
with $C=D=\mu/k^2$, $m=m_*$, $\alpha=\alpha_*$ and $A=A_wk^2$.
They therefore suffice. The root coordinates need not themselves be
bounded; only their displayed Gram coefficients must be bounded.

At $u,v$ infinitesimal use
\[
 E=(uW-Y)/k,\quad R=v(uX-Z)/k,\quad V=vW,\quad K=kuW,
 \quad X,Z.
\]
The exact amplitudes are
\[
 J=k(E+R),\quad M=D_uV-X-uZ-uvkE,\quad
 vY-Z=uV-Z-vkE,\quad kY=K-k^2E.
\]
Their full residue form is
\[
 4\mu\{E^2+(E+R)^2\}+4m_*(V-X)^2+4\alpha_*Z^2
       +4\overline{A_w}K^2+2X^2,                    \tag{UA13}
\]
which is strict. Literal zero testing angles are included without
any division by a testing angle.

For bounded $u$ and finite unit $v$, put
$E=(uW-Y)/k$, $F=J/k$. Then
\[
 Y=uW-kE,\quad Z=uX+k(E-F)/v.
\]
The roots $(E,F,W,X)$ have strict residue
\[
 4\mu D_vE^2+4\mu F^2+
 \{4m_*D_u^2+(4\alpha_*+2)D_uu^2\}(vW-X)^2
 +2D_vD_uX^2.                                      \tag{UA14}
\]
For bounded $u$ and $v=\infty$ write $q=1/v$ and homogenize by $q^2$.
The same roots are defined using the homogeneous $J$, and
\[
 Y=uW-kE,\qquad Z=uX+kqE-kF.
\]
The strict residue is (UA14) with $D_v=1$ and $(vW-X)^2$ replaced
by $W^2$. This includes the literal endpoint $q=0$.

For $u=\infty$, put $p=1/u$, $D_p=1+p^2$, and homogenize by $p^2$.
Write $r_p=W-pY$, $Q_p=X-pZ$, and $P=vY-Z$. Then
\[
 J=r_p+vQ_p,\qquad M=D_pP-pQ_p+pv r_p.
\]
For $v$ infinitesimal use the six roots
\[
 E=r_p/k,\quad F=vQ_p/k,\quad P,\quad R=pQ_p,
 \quad K=kY,\quad S=pZ.
\]
The signed core is exactly
\[
 2k^2(E+F)^2+2D_pP^2+2D_vR^2-4D_pD_vRP
    +4D_pD_vpFK+2D_pD_vS^2.
\]
Its cross-term reduction uses $RvY=pFK$, an exact identity.
Adding every capped packet gives the strict residue
\[
 4\mu\{E^2+(E+F)^2\}+(4m_*+2)(P-R)^2
 +4\alpha_*P^2+4\overline{A_w}K^2+2S^2.             \tag{UA15}
\]
For finite unit $v$, instead use $E=r_p/k$, $F=J/k$, $P,K,S$.
Since $Q_p=k(F-E)/v$, the strict residue is
\[
 4\mu D_vE^2+4\mu F^2+(4m_*+4\alpha_*+2)P^2
 +4\overline{A_w}D_vK^2+2D_vS^2.                   \tag{UA16}
\]
Finally, at both infinite endpoints put $q=1/v$ and homogenize by
$p^2q^2$. Use
\[
 E=(W-pY)/k,\quad F=(qW+X-pqY-pZ)/k,\quad
 P=Y-qZ,\quad K=kY,\quad S=pZ.
\]
Here $M=D_pP+pkD_qE-pqkF$ and the entire signed core is
\[
 2k^2F^2+2D_pP^2+2D_qp^2k^2(F-qE)^2
 +4D_qD_pk(F-qE)S+2D_qD_pS^2.
\]
The strict residue is (UA16) with $D_v=1$.

\paragraph{Exhaustion and restoration.}
A testing angle is infinitesimal, a bounded unit, or projectively
infinite. Equations (UA13)--(UA16) cover every pair except ordinary
$u$ with infinitesimal $v$. The local theorem covers the equal-residue
contact, and (UA4)--(UA12) cover its complement. Exact sign gauges are
\[
 (u,W,X)\mapsto(-u,-W,-X),\qquad
 (v,W,Y)\mapsto(-v,-W,-Y);
\]
all unspecified variables are fixed. Both preserve the full parent,
so negative ordinary contact residues are included. All endpoint
frames have been written after homogenization, including literal
projective endpoints. Each proof has an exact bounded Gram with
strict real residue; the strict-residue Gram lemma gives tensor SOS.
Restoring all retained packets completes the original parent.
Real and imaginary parts treat complex states. This closes the
entire strong C2 UI admission and its exact symmetry image HU.

\paragraph{Reproducibility.}
The recovered endpoint and allocation identities are replayed by
two scripts:
\begin{itemize}
\item \path{check_ui_projective_recovery.py}: seven exact checks.
\item \path{check_ui_all_tests_recovery.py}: thirty-one exact checks.
\end{itemize}
Those scripts verify symbolic equalities, sign gauges, all five
pole-free formulas, determinant factorization, contact derivatives,
and displayed endpoint residue forms. The explicit real interior
approximation also incorporates the same-system mathematical audit
in \path{ui_boundary_audit.tex}; it is not an independent
external verification. Positivity, boundedness and
exhaustion are supplied by the preceding analytical arguments;
passing algebraic checks alone is not used as their substitute.
